\subsection{An integrated nonsymmetric realization}

Section~\ref{sec:v44-realization} constructs a periodic family in which the
hypotheses of the intrinsic theorem hold simultaneously.  There are two
analytic obstacle types with disjoint curvature ranges, four prescribed
selected channels, two independent deck cycles based at one channel frame,
and a signature-rigid spanning tree.  Explicit support functions exclude
both proper and reflection symmetries.  Uniform separation and third-obstacle
clearance estimates hold for every lattice translate, while the lattice
shape, scale and the placement of the second obstacle vary.

For this family the registration step has a particularly transparent form.
Once the channel laws have yielded complete curve images, their support
coefficients of orders two and three determine relative rotations by a
single complex quotient.  Differences of the registered second-type
centers then give the two lattice vectors.  These support coefficients
are computed from recovered images; they are neither measured ordinates
nor additional calibration data.  Theorem~\ref{thm:v44-realized-family},
Proposition~\ref{prop:v44-harmonic-registration}, and
Corollary~\ref{cor:v44-eight-laws} give the construction, the registration
formula, and the complete observation-to-table reconstruction, respectively.
The fourth channel supplies a redundant holonomy identity in that same
frame.  Proposition~\ref{prop:v44-analytic-neighborhood} proves persistence
under independent analytic boundary perturbations, including higher
Fourier modes.  Thus the demonstration is not confined to fitting the
parameters of two trigonometric supports.

The explicit formulas concern geometric realization and the final matching
stage.  They do not replace the nonlinear relative limit or the smooth
finite-remainder contact inverse, and their registration stability is not
an effective analytic-continuation rate.  The construction keeps the
marked channels, signed conventions and prescribed gates of the original
observation map.  The richer planar-position benchmark remains distinct.
